\section{Related Work}\label{sec:related_work}
\xiangyu{simplify the description of each part in the related work because they are far from the focus of this paper.}
\subsection{Pure PTX injection for GPU Telemetry}
The Neutrino framework\cite{neutrino} pioneers zero-instrumentation GPU monitoring by leveraging PTX injection and hooks on CUDA runtime libraries. Without modifying application code, Neutrino attaches to user-level CUDA APIs to extract memory‐allocation patterns, kernel launch events, and resource‐usage metrics. It maps these events into isolated memory and delivers them to user space for visualization or offline analysis. This approach demonstrates that fine-grained GPU observability can be achieved with minimal overhead, and it provides a foundation upon which more advanced dynamic instrumentation and programmable tracing techniques can be built. However, it does not have the dynamism, CPU-GPU collaboration, and safety requirements that \sys provides.
   
\subsection{AI Observability Infrastructure}
Meta’s multi-layered observability stack\cite{ai-observability} combines low-overhead telemetry (via Dynolog and DCGM) with selective tracing (via CUPTI) to collect fleet-wide GPU metrics. Profiling tools such as PyTorch Profiler, Kineto, and Strobelight/BPF enable secure, real-time introspection of CUDA and PyTorch operations. An automated Analysis Platform stitches traces across hosts and highlights common inefficiencies—data starvation, kernel imbalances, or memory leaks—so developers can target optimizations effectively. Resource-attribution dashboards then allocate GPU hours or FLOPs by model or team, guiding long-term tuning efforts. Large-scale data centers require automated telemetry collection and analysis. Google’s GPU sharing on Kubernetes highlights orchestration at scale, while Meta’s Performance Profiling Platform automates trace stitching, anomaly detection, and trend analysis across thousands of jobs. These platforms underscore the necessity of coupling low-overhead data collection with intelligent, fleet-wide aggregation—an approach that informs \sys’s design for seamless integration into existing observability pipelines.

\subsection{Reducing Energy Usage in LLM training and inference}
Neurosurgeon\cite{neurosurgeon} introduces a framework for partitioning deep neural network inference between mobile devices and cloud servers to optimize end-to-end latency and energy consumption. By profiling layer-wise execution characteristics and network conditions, Neurosurgeon determines optimal split points in the model where computation should shift from edge to cloud. This collaborative intelligence approach enables resource-constrained devices to leverage cloud GPU power dynamically, reducing on-device latency and energy usage while maintaining accuracy. Perseus\cite{chung2024reducing} addresses the growing energy bottleneck in large-scale generative AI training by identifying and mitigating \emph{energy bloat}—energy consumed without contributing to throughput. It categorizes bloat into intrinsic (workload imbalance across pipeline parallel stages) and extrinsic (straggler-induced stalls in multi-pipeline training) sources. Perseus constructs the iteration time–energy Pareto frontier via an efficient graph–cut–based algorithm, then dynamically adjusts GPU frequency plans per computation to minimize total energy consumption without sacrificing throughput. Evaluations on models like GPT-3 demonstrate up to 30\% energy savings with negligible performance loss.
